\subsection{An explicit ELBO approximation and compatible sparsity control}
\label{sec:implementedelbo}
The automatic single-matrix interface now reports an explicit variational
diagnostic alongside its Monte Carlo EM loss. It also honors the learned
prior settings \path{penalty}, \path{n_epochs} and \path{batch_size}.
These additions do not replace the joint sampling kernel by a variational
E-step. The coupled \texttt{JointATACRNA} API retains its separate settings
and log-joint output.

\paragraph{A common regularized target.}
Write \(U\) for all augmented latent variables, including mixture labels
and HMM paths. For node \(u\) with spike probability
\(\pi_{0u}(P_{iu};\theta_u)\) and \(\lambda_u\geq1\), define
\begin{equation}
R_\lambda(U,\theta)=\sum_{i,u}(\lambda_u-1)
\log\pi_{0u}(P_{iu};\theta_u).
\label{eq:jointregularizer}
\end{equation}
Factor-axis priors contribute the analogous sum over feature rows. This is
the local learned-prior log-spike regularizer, evaluated at latent parents.
Since \(R_\lambda\leq0\), it has an auxiliary-observation interpretation:
conditionally independent events \(O_{iu}\) satisfy
\[
\Pr(O_{iu}=1\mid U,\theta)
=\pi_{0u}(P_{iu};\theta_u)^{\lambda_u-1},
\qquad O_{iu}=1\text{ is observed}.
\]
Consequently the regularized joint density is
\begin{equation}
p_\theta(Y,U,O=1)=p_\theta(Y,U)\exp\{R_\lambda(U,\theta)\}.
\label{eq:auxiliarytarget}
\end{equation}
It is a slice of a normalized augmented probability model, not a silently
renormalized conditional mixture. Conditional on \(Y,O=1\), its normalizing
constant is the corresponding evidence. Parameter learning by complete-data
EM therefore does not omit an additional parameter-dependent normalizer.
All original conditional mixture weights still sum to one. These events
express the chosen regularizer; they are not additional measured molecular
data. \(\lambda=1\) recovers the original model. Unspiked EMDN requires
\(\lambda=1\); spiked EMDN supports this regularization.

For a candidate parent \(v\), the correction in
\eqref{eq:algorithmcorrection} must replace \(D_u(v)\) by
\begin{equation}
D_u^\lambda(v)=\sum_{c\in\ch(u)}\left[
\log g_{\theta c}(V_c,Z_c\mid P_c[v])
+(\lambda_c-1)\log\pi_{0c}(P_c[v];\theta_c)\right].
\label{eq:regularizedchild}
\end{equation}
The node's own regularizer is constant in its value because its parents are
fixed during that move. A child's regularizer is generally not constant in
the candidate parent. The implementation includes both the full child
density and its regularizer; omitting the latter would decouple the
sampling and learning objectives.

The fixed-draw parameter loss becomes
\begin{equation}
J^\lambda_{u,t}(\theta_u)=-\frac{1}{SN}\sum_{s,i}
\left[\log g_{\theta u}(V_{iu}^{(s)},Z_{iu}^{(s)}\mid P_{iu}^{(s)})
+(\lambda_u-1)\log\pi_{0u}(P_{iu}^{(s)};\theta_u)\right].
\end{equation}
Mini-batches contain intact parent/value/label pairs and average both terms
on the same scale. One epoch visits the entire saved dataset. After each
epoch, the full-dataset loss is evaluated and the best finite state is
retained, including the initial state. Thus batch size changes optimization
noise and the number of gradient steps, not regularization strength.
For an intercept-only gate with a saved spike fraction \(r_0\), minimizing
this loss gives
\(\widehat\pi_0=(r_0+\lambda-1)/\lambda\), explicitly exhibiting its
shrinkage toward the spike. This is also an independent numerical test.

\paragraph{A normalized posterior approximation.}
For scalar nodes outside an HMM, fit a labelled mixture to saved states:
\[
q_{iu}(dv,z)=\widehat r_{iuz}\,
\begin{cases}
\delta_{a_z}(dv),&z\text{ is an atomic component},\\
\N(dv;\widehat m_{iuz},\widehat v_{iuz}),&z\text{ is continuous}.
\end{cases}
\]
The label probabilities are empirical frequencies. Continuous means and
variances are fitted within labels, with a positive variance floor; fewer
than two visits use a reference normal-means posterior variance. Unvisited
labels have zero \(q\)-mass. Fixed atoms retain their exact locations.
The product approximation factorizes across axes and factor columns, and
across rows for scalar priors. This is a diagnostic projection of the
joint sample, not the joint posterior used to reconstruct the signal.

An HMM column instead uses the normalized full-path distribution
\[
q_k(dv,h,z)\propto p_{H,k}(dv,h,z)
\exp\left\{\sum_j(-a_jv_j^2/2+b_jv_j)\right\},
\]
where the reference statistics are fixed using the saved mean factor
matrices. Its normalizer is obtained by the forward recursion, and draws
use whole-path FFBS followed by component/value draws. Forbidden transitions,
nonzero atoms and one-sided supports remain intact. The density is available
on the same augmented reference measure as the model.

\paragraph{The reported bounds and numerical estimates.}
For this explicit normalized \(q\), define
\begin{align}
\LL(q,\theta)&=\E_q[\log p_\theta(Y,U)-\log q(U)]
\leq\log p_\theta(Y),\\
\LL_\lambda(q,\theta)&=\LL(q,\theta)+\E_qR_\lambda(U,\theta)
\leq\log p_\theta(Y,O=1).
\label{eq:diagnosticelbo}
\end{align}
New independent draws \(\widetilde U^{(d)}\sim q\), using a separate random
generator, estimate the expectations by averages of \(\log p-\log q\),
with the regularizer added for \(\LL_\lambda\). The categorical and
continuous entropy contributions are both included. A discrete empirical
entropy of continuous MCMC draws would not give this bound.

The sample standard deviation of each log-ratio divided by the square root
of the evaluation sample size estimates its Monte Carlo standard error.
It is conditional on fitted \(q,\theta\), and does not quantify mixing,
hyperparameter uncertainty or the variational gap. Although the expectations
in \eqref{eq:diagnosticelbo} are lower bounds, a finite Monte Carlo estimate
can exceed the evidence. No monotonicity is claimed: Monte Carlo EM trains
on joint states, whereas this explicit marginal approximation is used for
diagnosis. Refitting \(q\) between rounds also changes the evaluated bound.

In the automatic interface, \path{history_obj} stores
\(-\widehat\LL_\lambda\), so lower is better. Each \path{history_elbo}
record contains the ordinary and regularized estimates, their standard
errors, and the evaluation phase. The final estimate is
\path{result.elbo_estimate}; \path{elbo_draws} controls evaluation effort.
The sampled augmented log joint is retained separately in
\path{result.joint_posterior.log_joint}. Numerical checks compare this
construction with exact finite-state ELBOs, analytic Gaussian moments,
mixed-distribution entropy, and HMM normalization by quadrature and path
enumeration.
